% =======================================================
\section{Introduction}
\label{sec:intro}
% =======================================================

Reconstructing a computed tomography (CT) image from few projection angles or at low photon counts is an
ill-posed inverse problem. From measurements
\begin{equation}
y = Ax + n,
\end{equation}
where $A$ is the forward (Radon) operator and $n$ the noise, the image $x$ cannot be recovered stably when
$A$ is rank-deficient or ill-conditioned, as it is for sparse-view data, without prior knowledge about $x$.
The classical remedy is explicit regularization,
\begin{equation}
\hat{x} = \arg\min_x \; \tfrac{1}{2}\lVert Ax - y \rVert_2^2 + \lambda R(x),
\end{equation}
with a hand-crafted penalty $R$ such as Tikhonov or total variation (TV) \citep{rudin1992tv}.

Learned priors replace $R$ by statistics estimated from data. End-to-end networks, including unrolled
iterations, are trained for a given operator and noise level \citep{jin2017,adler2018}. Denoiser-based methods,
Plug-and-Play (PnP) \citep{venkatakrishnan2013plugandplay} and regularization by denoising \citep{romano2017}
(reviewed in \citep{kamilov2023}), instead insert a trained denoiser such as DnCNN \citep{zhang2017dncnn}
in place of the proximal operator of $R$ in an iterative scheme. Learned methods can be much more accurate than
classical reconstruction on large benchmarks: in the LoDoPaB-CT dataset paper, FBP followed by a U-Net improves the PSNR by several dB over FBP \citep{leuschner2021}. We study PnP with a DnCNN prior because
the prior is trained once, independently of the operator, so the same prior can be placed in the same
data-fidelity iteration as the classical baselines and tested when the geometry, the dose or the image
distribution changes. We make no claim about end-to-end networks trained for a specific task, and our
conclusions concern one small denoiser architecture.

Reconstruction methods are usually compared by PSNR and SSIM on test images drawn from the distribution of the
training set. Several concerns are known to be left open by such comparisons. Learned reconstructions can be
unstable under small perturbations \citep{antun2020instabilities} and sensitive to distribution shift
in accelerated MRI \citep{darestani2021shift}; they can contain plausible structure that is not in the object
\citep{bhadra2021}; and denoiser-based iterations with an unconstrained CNN denoiser come without convergence
guarantees \citep{ryu2019pnp,chan2016pnp}. Moreover, image quality for a diagnostic purpose is better measured
by performance at a task than by a pixel-wise error \citep{barrett1993,barrett2004}. Finally, the verdict depends
on the baseline: a regularizer whose weight and iteration count are not tuned for each condition is a weak
comparison.

This paper is an evaluation study, not a new method. It asks three questions about a PnP-DnCNN prior compared with
tuned FBP, Tikhonov and TV reconstruction:
\begin{enumerate}
\item[\textbf{Q1}] Does the learned prior keep its advantage over TV when the image distribution changes?
\item[\textbf{Q2}] Can we find cases where it creates lesion-like structure that the ground truth does not
support, as measured by a location-known local-contrast statistic?
\item[\textbf{Q3}] Do the conclusions hold over a range of noise levels and reconstruction conditions?
\end{enumerate}
We answer them at two scales. In a controlled setting of $28\times28$ digit images (MNIST-derived), noise level,
number of angles, forward-model error, denoiser seeds and the training distribution of the prior can all be varied
and repeated; the same questions are then examined with a matrix-free operator and Poisson noise. We then ask which
conclusions survive on real $362\times362$ slices of LoDoPaB-CT with a patient-disjoint split, three angle counts
and three dose levels. In every experiment the data are split into disjoint training, validation and test sets,
each method is tuned on the validation set for each condition, and differences are paired (Section~\ref{sec:method}).
Besides PSNR and SSIM, we measure the detectability of synthetic lesions inserted into the images (paired AUC,
lesion retention and false lesion-like blobs). The first digit experiments (E1--E5) use inverse-crime data; later
experiments simulate the measurements on a finer grid \citep{wirgin2004}, and
Section~\ref{sec:crime-results} quantifies the difference.

In brief: on digits a matched prior beats TV by about $3.4$\,dB at moderate noise, and a mismatched prior loses to
it by $3.1$--$4.6$\,dB. On real CT, one PnP-DnCNN prior is roughly on par with a well-tuned TV in PSNR and slightly
below it in SSIM, yet it detects small high-contrast lesions better than that TV (significantly in six of nine
conditions), which itself detects them worse than plain FBP.

Our contributions are:
\begin{itemize}
\item An evaluation protocol for comparing a PnP prior with classical regularization: per-condition tuning of
every method, a TV control run for twice the iterations (TV-long), paired statistics, and data generation on a
finer grid with a sensitivity analysis (Section~\ref{sec:crime-results}).
\item Robustness evidence on digit images: forward-model mismatch, distribution shift with matched-prior and seed
controls, the effect of the iteration count, and an error decomposition along the singular vectors of the forward
operator (Sections~\ref{sec:e3}--\ref{sec:seeds-results}).
\item A lesion-insertion evaluation with an operating-point study of whether re-tuning can restore detectability
(Sections~\ref{sec:lesion-results} and \ref{sec:oppoint-results}).
\item A matrix-free operator with exact adjoint, validated against the dense one, and a Poisson-noise study
(Sections~\ref{sec:mf-results}--\ref{sec:lesion-poisson}).
\item A patient-disjoint study on LoDoPaB-CT over three angle counts and three doses, with a simulator check and a
shift to non-anatomical phantoms (Sections~\ref{sec:e14}--\ref{sec:e19}). The study reports negative and mixed
results, an early comparison that we discarded, and the analysis steps taken after seeing results
(Section~\ref{sec:history3}); nothing was pre-specified.
\end{itemize}

Section~\ref{sec:related} reviews related work, Section~\ref{sec:method} the methods, Section~\ref{sec:experiments}
reports experiments E1--E19, and Section~\ref{sec:discussion} answers the three questions.